% !PRINTED 333
\weberpage{333}
%<<CONT
dass man den zwischen den beiden Curvenpunkten $\alpha, f(\alpha)$ und
$\alpha + \delta, f(\alpha + \delta)$ verlaufenden Curvenbogen durch die Sehne
ersetzt. Berücksichtigt man auch die zweite Differenz, so wird der durch die
drei auf einander folgenden Punkte mit den Abscissen $\alpha, \alpha + \delta,
\alpha + 2\delta$ gehende Curvenbogen durch den Bogen einer Parabel ersetzt,
die durch dieselben Punkte geht und deren Axe der Ordinatenaxe parallel ist;
und diese Parabel wird sich der Curve noch enger anschliessen als die Sehne.
Je kleiner das Intervall $\delta$ ist, um so weniger werden die höheren
Differenzen von Einfluss sein.

Wie weit man also bei der Annäherung zu gehen hat, das hängt nicht nur von der
Genauigkeit ab, mit der man $f(x)$ zu kennen wünscht, sondern wesentlich auch
von der Dichtigkeit der Werthe $\alpha, \alpha + \delta, \alpha + 2\delta
\ldots$, für die die Function bekannt ist. Auf diesen Sätzen beruht die
Einrichtung unserer Tabellenwerke, besonders der Logarithmentafeln. Es
handelt sich dabei freilich nicht um ganze rationale Functionen; allein bei
den stetigen Functionen überhaupt gelten hier dieselben Gesetze. Man findet
in den Tafeln daher auch neben den Werthen $f(\alpha), f(\alpha + \delta),
f(\alpha + 2\delta) \ldots$ die Werthe der ersten oder der beiden ersten
Differenzen angegeben. Bei den gebräuchlichen siebenstelligen Tafeln genügt
die erste Differenz. In der zehnstelligen Tafel \glqq Thesaurus
logarithmorum\grqq{} von \gesperrt{Vega} sind auch die zweiten Differenzen
angegeben und müssen bei ganz scharfen Rechnungen berücksichtigt werden.

Unsere Interpolationsformeln lassen sich mit Nutzen anwenden, um die Wurzeln
der Gleichungen zu berechnen, oder, genauer ausgedrückt, die auf anderem
Wege gefundenen Näherungswerthe zu verbessern.

Wir können die Aufgabe so formuliren, dass zu einem gegebenen, zwischen
$f(\alpha)$ und $f(\alpha + \delta)$ gelegenen Werth von $f(x)$ der
zugehörige Werth von $x$ gefunden werden soll.

Wir betrachten $x = \alpha$ als einen ersten Näherungswerth. Setzen wir nun
\[
f(x) - f(\alpha) = \varDelta,
\]
so ist $\varDelta$ ein gegebener Werth von demselben Zeichen wie
$\varDelta_\alpha = f(\alpha + \delta) - f(\alpha)$ und absolut kleiner als
$\varDelta_\alpha$. Setzen wir noch $x - \alpha = u$, so giebt die Formel (4)
\weq{6}{\varDelta = \frac{\varDelta_\alpha u}{\delta} + \frac{\varDelta'_\alpha}{\delta^2}\,\frac{u(u-\delta)}{2} + \cdots}
